\documentclass[10pt,a4paper]{amsart}
\usepackage{amsmath,amssymb,amsthm,graphicx}
% Retained minimal preamble for the Stacks Project chapter more-algebra.tex (More on Algebra).
% The chapter file pinned-original-source.tex carries only \input{preamble} on its first line, so no
% preamble text lives inside the chapter itself.  The Stacks Project uses one single shared preamble for
% every chapter; the snapshot stages it identically as sources/stacks-chow/preamble.tex and
% sources/stacks-etale-cohomology/preamble.tex.  The block below is copied mechanically, line for line,
% from sources/stacks-chow/preamble.tex lines 14-18 (the Xy-pic setup) and 154-217 (the native theorem
% declarations together with the mathematical macros).  Nothing was added, removed or redefined; the
% wrapper class and the amsmath/amssymb/amsthm/graphicx base packages are supplied by the assembler.
% Chapter: more-algebra.tex -- More on Algebra.
% For commutative diagrams we use Xy-pic
\usepackage[all]{xy}

% We use 2cell for 2-commutative diagrams.
\xyoption{2cell}

% Theorem environments.
%
\theoremstyle{plain}
\newtheorem{theorem}[subsection]{Theorem}
\newtheorem{proposition}[subsection]{Proposition}
\newtheorem{lemma}[subsection]{Lemma}

\theoremstyle{definition}
\newtheorem{definition}[subsection]{Definition}
\newtheorem{example}[subsection]{Example}
\newtheorem{exercise}[subsection]{Exercise}
\newtheorem{situation}[subsection]{Situation}

\theoremstyle{remark}
\newtheorem{remark}[subsection]{Remark}
\newtheorem{remarks}[subsection]{Remarks}

\numberwithin{equation}{subsection}

% Macros
%
\def\lim{\mathop{\mathrm{lim}}\nolimits}
\def\colim{\mathop{\mathrm{colim}}\nolimits}
\def\Spec{\mathop{\mathrm{Spec}}}
\def\Hom{\mathop{\mathrm{Hom}}\nolimits}
\def\Ext{\mathop{\mathrm{Ext}}\nolimits}
\def\SheafHom{\mathop{\mathcal{H}\!\mathit{om}}\nolimits}
\def\SheafExt{\mathop{\mathcal{E}\!\mathit{xt}}\nolimits}
\def\Sch{\mathit{Sch}}
\def\Mor{\mathop{\mathrm{Mor}}\nolimits}
\def\Ob{\mathop{\mathrm{Ob}}\nolimits}
\def\Sh{\mathop{\mathit{Sh}}\nolimits}
\def\NL{\mathop{N\!L}\nolimits}
\def\CH{\mathop{\mathrm{CH}}\nolimits}
\def\proetale{{pro\text{-}\acute{e}tale}}
\def\etale{{\acute{e}tale}}
\def\QCoh{\mathit{QCoh}}
\def\Ker{\mathop{\mathrm{Ker}}}
\def\Im{\mathop{\mathrm{Im}}}
\def\Coker{\mathop{\mathrm{Coker}}}
\def\Coim{\mathop{\mathrm{Coim}}}

% Boxtimes
%
\DeclareMathSymbol{\boxtimes}{\mathbin}{AMSa}{"02}

%
% Macros for moduli stacks/spaces
%
\def\QCohstack{\mathcal{QC}\!\mathit{oh}}
\def\Cohstack{\mathcal{C}\!\mathit{oh}}
\def\Spacesstack{\mathcal{S}\!\mathit{paces}}
\def\Quotfunctor{\mathrm{Quot}}
\def\Hilbfunctor{\mathrm{Hilb}}
\def\Curvesstack{\mathcal{C}\!\mathit{urves}}
\def\Polarizedstack{\mathcal{P}\!\mathit{olarized}}
\def\Complexesstack{\mathcal{C}\!\mathit{omplexes}}
% \Pic is the operator that assigns to X its picard group, usage \Pic(X)
% \Picardstack_{X/B} denotes the Picard stack of X over B
% \Picardfunctor_{X/B} denotes the Picard functor of X over B
\def\Pic{\mathop{\mathrm{Pic}}\nolimits}
\def\Picardstack{\mathcal{P}\!\mathit{ic}}
\def\Picardfunctor{\mathrm{Pic}}
\def\Deformationcategory{\mathcal{D}\!\mathit{ef}}

% --- appended 2026-10-12 (Chief): the project own definitions of the reference/slogan/history
% display environments, copied mechanically from the shared Stacks preamble lines 8-12
% (\usepackage{verbatim} plus the three \newenvironment lines, which the chapter preamble keeps
% verbatim from the project preamble). Nothing added, removed or redefined: these environments
% carry no mathematical content and the project itself typesets them as comments. Without these
% three lines a counted statement whose own text uses them fails with Environment slogan undefined.
\usepackage{verbatim}
\newenvironment{reference}{\comment}{\endcomment}
\newenvironment{slogan}{\comment}{\endcomment}
\newenvironment{history}{\comment}{\endcomment}
\usepackage{cleveref}
\crefname{originaldefinition}{original definition}{original definitions}
\Crefname{originaldefinition}{Original definition}{Original definitions}
\makeatletter
\newlabel{definition-extension-discrete-valuation-rings}{{09E4}{source}{}{sourcepointer-1}{}}
\newlabel{definition-extension-discrete-valuation-rings@cref}{{[originaldefinition][0][]09E4}{[1][0][]source}{}{}{}}
\makeatother
\crefname{originallemma}{original lemma}{original lemmas}
\Crefname{originallemma}{Original lemma}{Original lemmas}
\makeatletter
\newlabel{algebra-lemma-characterize-dvr}{{00PD}{source}{}{sourcepointer-2}{}}
\newlabel{algebra-lemma-characterize-dvr@cref}{{[originallemma][0][]00PD}{[1][0][]source}{}{}{}}
\makeatother
\crefname{originallemma}{original lemma}{original lemmas}
\Crefname{originallemma}{Original lemma}{Original lemmas}
\makeatletter
\newlabel{algebra-lemma-minimal-over-1}{{00KV}{source}{}{sourcepointer-3}{}}
\newlabel{algebra-lemma-minimal-over-1@cref}{{[originallemma][0][]00KV}{[1][0][]source}{}{}{}}
\makeatother
\crefname{originallemma}{original lemma}{original lemmas}
\Crefname{originallemma}{Original lemma}{Original lemmas}
\makeatletter
\newlabel{algebra-lemma-normal-domain-intersection-localizations-height-1}{{031T}{source}{}{sourcepointer-4}{}}
\newlabel{algebra-lemma-normal-domain-intersection-localizations-height-1@cref}{{[originallemma][0][]031T}{[1][0][]source}{}{}{}}
\makeatother
\crefname{originallemma}{original lemma}{original lemmas}
\Crefname{originallemma}{Original lemma}{Original lemmas}
\makeatletter
\newlabel{algebra-lemma-finite-projective-descends}{{00O1}{source}{}{sourcepointer-5}{}}
\newlabel{algebra-lemma-finite-projective-descends@cref}{{[originallemma][0][]00O1}{[1][0][]source}{}{}{}}
\makeatother
\title{A root of a function in a weakly unramified extension of normal local domains descends to the smaller ring}
\author{The Stacks Project Authors}
\date{stacks-project:more-algebra.tex:local-snapshot-20261012; GNU FDL 1.2 or later}
\begin{document}
\maketitle
\setcounter{subsection}{1}
\noindent\textit{Editorial note.} All delivered mathematical text is the author's own, copied from the pinned Stacks Project chapter file more-algebra.tex; no statement or proof step was rewritten, reordered or removed. The counted span is the author environment and its own proof environment. References inside the span are handed to the assembler as navigation stubs that carry the canonical tag of the cited tag only; whether each of them is load-bearing is recorded on the stub itself.\par\medskip
\noindent\textit{Allowed standard prerequisites (editorial).}\begin{itemize}
\item The local rings at height one primes of a Noetherian normal domain are discrete valuation rings, and the notion of a weakly unramified extension of such rings is the one defined for extensions of discrete valuation rings. used at proof lines 37401-37404. Makes the hypothesis of the statement meaningful and supplies the valuations used in every later step.
\item Primes minimal over one element of a Noetherian domain have height one. invoked at proof lines 37405-37407. Shows that the primes minimal over the function are exactly the height one primes whose valuations the proof computes.
\item An element of a normal domain whose valuation at every height one prime is bounded below by that of another element lies in the multiple of the second element. invoked at proof lines 37430-37436 and again at proof lines 37446-37448 and at the end at proof lines 37473-37476. Underlies the claim that the extension of the ideal cuts out exactly the elements with the prescribed valuation bounds, and the final comparison of the two generators.
\item A module over the smaller ring which becomes free of rank one after a faithfully flat base change is itself free of rank one. invoked at proof lines 37468-37470. Descends the rank one free module property from the larger ring to the smaller one, producing the element whose power differs from the given function by a unit.
\end{itemize}\medskip
\ifdefined\hypertarget\hypertarget{sourcepointer-1}{}\fi
\noindent\textit{Source reference.} Original definition 09E4: sources/stacks-more-algebra/pinned-original-source.tex, source lines 32912--32912. The proof recalls at proof lines 37403-37404 that the hypothesis is the notion of that definition; the definition is a pointer fixing terminology, and every valuation computation using it is inside the counted span.\par\medskip
\ifdefined\hypertarget\hypertarget{sourcepointer-2}{}\fi
\noindent\textit{Source reference.} Original lemma 00PD: sources/stacks-chow/tags.txt, source lines 1--1. Used at proof lines 37401-37403 to see that the local rings in question are discrete valuation rings; the valuations taken with respect to them are computed inside the counted span.\par\medskip
\ifdefined\hypertarget\hypertarget{sourcepointer-3}{}\fi
\noindent\textit{Source reference.} Original lemma 00KV: sources/stacks-chow/tags.txt, source lines 1--1. Invoked at proof lines 37405-37407 to see that the primes minimal over f are the height one primes; the enumeration of the primes above each of them is inside the counted span.\par\medskip
\ifdefined\hypertarget\hypertarget{sourcepointer-4}{}\fi
\noindent\textit{Source reference.} Original lemma 031T: sources/stacks-chow/tags.txt, source lines 1--1. Used three times inside the argument, at proof lines 37430-37436 for the injectivity claim, at proof lines 37446-37448 for the description of the extended ideal and at the end at proof lines 37473-37476 for the comparison of the two generators; each time the element whose membership is being tested is produced inside the counted span.\par\medskip
\ifdefined\hypertarget\hypertarget{sourcepointer-5}{}\fi
\noindent\textit{Source reference.} Original lemma 00O1: sources/stacks-chow/tags.txt, source lines 1--1. Invoked at proof lines 37468-37470 to descend the property of being free of rank one from the larger ring to the smaller one; the identification of the extended ideal with a rank one free module is inside the counted span.\par\medskip
\begin{lemma}
\label{lemma-extension-normal-domains-and-roots}
Let $A \to B$ be a flat local homomorphism of Noetherian local normal domains.
Let $f \in A$ and $h \in B$ such that $f = w h^n$ for some $n > 1$ and some
unit $w$ of $B$. Assume that for every height $1$ prime
$\mathfrak p \subset A$ there is a height $1$ prime
$\mathfrak q \subset B$ lying over $\mathfrak p$
such that the extension $A_\mathfrak p \subset B_\mathfrak q$ is
weakly unramified. Then $f = u g^n$ for some $g \in A$ and unit $u$ of $A$.
\end{lemma}

\begin{proof}
The local rings of $A$ and $B$ at height $1$ primes are
discrete valuation rings (Algebra, Lemma \ref{algebra-lemma-characterize-dvr}).
Thus the assumption makes sense (via
Definition \ref{definition-extension-discrete-valuation-rings}).
Let $\mathfrak p_1, \ldots, \mathfrak p_r$ be the primes of $A$ minimal
over $f$. These have height $1$ by
Algebra, Lemma \ref{algebra-lemma-minimal-over-1}.
For each $i$ let $\mathfrak q_{i, j} \subset B$, $j = 1, \ldots, r_i$
be the height $1$ primes of $B$ lying over $\mathfrak p_i$.
Say we number them so that  $A_{\mathfrak p_i} \to B_{\mathfrak q_{i, 1}}$
is weakly unramified.
Since $f$ maps to an $n$th power times a unit in $B_{\mathfrak q_{i, 1}}$
we see that the valuation $v_i$ of $f$ in $A_{\mathfrak p_i}$ is
divisible by $n$. Say $v_i = n w_i$ for some $w_i \geq 0$.
Consider the exact sequence
$$
0 \to I \to A \to
\prod\nolimits_{i = 1, \ldots, r}
A_{\mathfrak p_i}/\mathfrak p_i^{w_i}A_{\mathfrak p_i}
$$
defining the ideal $I$. Applying the exact functor $- \otimes_A B$ we obtain
an exact sequence
$$
0 \to I \otimes_A B \to B \to
\prod\nolimits_{i = 1, \ldots, r}
(A_{\mathfrak p_i}/\mathfrak p_i^{w_i}A_{\mathfrak p_i}) \otimes_A B
$$
Fix $i$. We claim that the canonical map
$$
(A_{\mathfrak p_i}/\mathfrak p_i^{w_i}A_{\mathfrak p_i}) \otimes_A B
\to
\prod\nolimits_{j = 1, \ldots, r_i}
B_{\mathfrak q_{i, j}}/\mathfrak q_{i, j}^{e_{i, j}w_i}B_{\mathfrak q_{i, j}}
$$
is injective. Here $e_{i, j}$ is the ramification index of 
$A_{\mathfrak p_i} \to B_{\mathfrak q_{i, j}}$.
The claim asserts that $\mathfrak p_i^{w_i}B_{\mathfrak p_i}$
is equal to the set of elements $b$ of $B_{\mathfrak p_i}$
whose valuation at $\mathfrak q_{i, j}$ is $\geq e_{i, j}w_i$.
Choose a generator $a \in A_{\mathfrak p_i}$ of the principal
ideal $\mathfrak p_i^{w_i}$. Then the valuation of $a$ at
$\mathfrak q_{i, j}$ is equal to $e_{i, j}w_i$. Hence,
as $B_{\mathfrak p_i}$ is a normal domain whose height one
primes are the primes $\mathfrak q_{i, j}$, $j = 1, \ldots, r_i$,
we see that, for $b$ as above, we have $b/a \in B_{\mathfrak p_i}$ by
Algebra, Lemma
\ref{algebra-lemma-normal-domain-intersection-localizations-height-1}.
Thus the claim.

\medskip\noindent
The claim combined with the second exact sequence above
determines an exact sequence
$$
0 \to I \otimes_A B \to B \to
\prod\nolimits_{i = 1, \ldots, r}
\prod\nolimits_{j = 1, \ldots, r_i}
B_{\mathfrak q_{i, j}}/\mathfrak q_{i, j}^{e_{i, j}w_i}B_{\mathfrak q_{i, j}}
$$
It follows that $I \otimes_A B$ is the set of elements $h'$ of $B$
which have valuation $\geq e_{i, j}w_i$ at $\mathfrak q_{i, j}$.
Since $f = wh^n$ in $B$ we see that $h$ has valuation
$e_{i, j}w_i$ at $\mathfrak q_{i, j}$. Thus $h'/h \in B$
by Algebra, Lemma
\ref{algebra-lemma-normal-domain-intersection-localizations-height-1}.
It follows that $I \otimes_A B$ is a free $B$-module of rank $1$
(generated by $h$).
Therefore $I$ is a free $A$-module of rank $1$, see
Algebra, Lemma \ref{algebra-lemma-finite-projective-descends}.
Let $g \in I$ be a generator. Then we see that
$g$ and $h$ differ by a unit in $B$. Working backwards we
conclude that the valuation of $g$ in $A_{\mathfrak p_i}$ is
$w_i = v_i/n$. Hence $g^n$ and $f$ differ by a unit in $A$
(by Algebra, Lemma
\ref{algebra-lemma-normal-domain-intersection-localizations-height-1})
as desired.
\end{proof}

\end{document}
